\subsection{Практические задания: базовые операции}
Этот блок опирается только на уже изученные основы: создание объекта \verb|Ledbar|, модель RGB, индексацию \verb|0..24| и вызов \verb|leds:set(...)|. Здесь не нужны таймеры, обработчик \verb|callback(event)| и конечные автоматы.

\begin{taskbox}[Одновременное включение всей светодиодной линейки зелёным цветом]
\indent Требуется реализовать алгоритм, обеспечивающий однократную установку зелёного цвета с компонентами \verb|{0,1,0}| на всех светодиодах линейки с индексами от \verb|0| до \verb|24|. Для этого необходимо создать объект светодиодной линейки, реализовать функцию массовой установки цвета и выполнить её вызов.

\noindent\textbf{Ожидаемый результат:} вся линейка равномерно светится зелёным цветом.
\end{taskbox}

\begin{taskbox}[Чередование цветов по чётности индекса]
\indent Требуется реализовать статический шаблон индикации, при котором светодиоды с чётными индексами светятся красным цветом \verb|{1,0,0}|, а с нечётными --- синим цветом \verb|{0,0,1}|. Изображение должно устанавливаться однократно, без анимации и таймеров.

\noindent\textbf{Ожидаемый результат:} на линейке виден устойчивый чередующийся цветовой узор.
\end{taskbox}

\begin{taskbox}[Индикатор заполнения]
\indent Необходимо реализовать однократное заполнение линейки от младших индексов к старшим: светодиоды с индексами от \verb|0| до некоторого значения \verb|k| включаются зелёным цветом \verb|{0,1,0}|, а остальные остаются выключенными. Значение \verb|k| можно задать прямо в программе.

\noindent\textbf{Ожидаемый результат:} на линейке отображается статическая «полоса загрузки» длиной \verb|k+1| светодиодов.
\end{taskbox}
